\documentclass[slidestop,compress,mathserif]{beamer}
%\documentclass[slidestop,compress,mathserif,handout]{beamer}

%\documentclass[xcolor=dvipsnames,handout]{beamer}
%\documentclass[xcolor=dvipsnames]{beamer}

%\documentclass[handout]{beamer}

%%% To get rid of solutions on handouts:
\newcommand{\soln}[1]{\textit{\textcolor{darkGray}{#1}}}				% For slides
%\newcommand{\soln}[1]{ }	% For handouts

% to get pausing to work properly on slides
\newcommand{\hide}[1]{#1}	% For slides
%\newcommand{\hide}[1]{ }	% For handouts
\usepackage{tikz}


%\usepackage{multicol}
\usepackage{amsfonts}
%\usepackage[pdftex,dvipsnames]{color}
\usepackage{graphicx}
\usepackage{subfigure}
%\usepackage{picinpar}
\usepackage{pifont}
\usepackage{pgf,pgfarrows,pgfnodes}
%\usepackage{wasysym,manfnt,phaistos,empheq}
\usepackage[english]{babel}
\usepackage{pgfpages}
\usepackage{natbib}
\usepackage{hyperref}
\usepackage{multimedia}
%\usepackage{amsfonts,amstext,amssymb,amsbsy,amsopn,amsthm,eucal,latexsym,mathrsfs}
\usepackage{amsmath,amsfonts,amstext,amssymb,amsbsy,amsopn,amsthm,eucal,latexsym,mathrsfs}
\usepackage{ulem}
\usepackage{setspace}
\usepackage{array}
%\usepackage{rotating}
\usepackage{multirow}
\usepackage{verbatim}
\usepackage{multicol}

\setbeamertemplate{navigation symbols}{}

%\usepackage{tikz}
%\usetikzlibrary{arrows,shapes,trees,backgrounds}


%\setbeameroption{show notes on second screen}
%\setbeameroption{show notes}
%\setbeameroption{show only notes}

\definecolor{links}{HTML}{2A1B81}
\hypersetup{colorlinks,linkcolor=,urlcolor=links}

\newtheorem*{principle}{Inscrutibility Principle}
\newtheorem*{punchline}{Punch Line}
\newtheorem{defn}{Definition}

\definecolor{Scarlet}{RGB}{140,17,17}
\definecolor{VassarRed}{RGB}{128,0,0}

% "dinglist" environment
  \renewenvironment{dinglist}[2][blue]
  {\begin{list}{\textcolor{blue}{\ding{#2}}}{}}{\end{list}}
  % Symbol definitions for these lists
  \newcommand{\DingListSymbolA}{43}
  \newcommand{\DingListSymbolB}{243}
  \newcommand{\DingListSymbolC}{224}
  \newcommand{\DingListSymbolD}{219}
  \newcommand{\DingListSymbolCheck}{52}
  \newcommand{\DingListSymbolCross}{56}


  \newenvironment{ballotenv}
{\only{%
\setbeamertemplate{itemize item}{\ding{45}}%
\setbeamertemplate{itemize subitem}{\ding{46}}%
\setbeamertemplate{itemize subsubitem}{\ding{46}}}} {}
\setbeamertemplate{itemize item}{\ding{49}}
\setbeamertemplate{itemize subitem}{\ding{47}}
\setbeamertemplate{itemize subsubitem}{\ding{47}}


%User defined colors: See colors section
\xdefinecolor{oiBlue}{rgb}{0.15, 0.35, 0.55}
\xdefinecolor{gray}{rgb}{0.5, 0.5, 0.5}
\xdefinecolor{darkGray}{rgb}{0.3, 0.3, 0.3}
\xdefinecolor{darkerGray}{rgb}{0.2, 0.2, 0.2}
\xdefinecolor{rubineRed}{rgb}{0.89,0,0.30}
\xdefinecolor{linkCol}{rgb}{0.11,0.49,0.95}	
\xdefinecolor{irishGreen}{rgb}{0,0.60,0}	
\xdefinecolor{darkturquoise}{rgb}{0.44, 0.58, 0.86}
\definecolor{lightGreen}{rgb}{0.533,0.765,0.42}
\xdefinecolor{Regalia}{HTML}{522D80}
\xdefinecolor{ClemsonOrange}{HTML}{EA6A20}

\definecolor{duke@LightGrey}{RGB}{200,200,200}\definecolor{DarkGreen}{RGB}{0,100,0}
\definecolor{Oranges}{RGB}{255,127,0}
\definecolor{LightGray}{RGB}{211,211,211}

%\setbeamertemplate{footline}{%
%  \raisebox{5pt}{\makebox[\paperwidth]{\hfill\makebox[10pt]{\scriptsize\insertframenumber}}}}

\setbeamercolor{equation background}{fg=black,bg=duke@LightGrey}
  % Boxed equation
  \newcommand{\eqbox}[2][0.6]{%
  \centerline{
  \begin{beamerboxesrounded}[lower=equation background,width=#1\hsize,shadow=true]{}
\parbox{#1\hsize}{%
      \[
        \textcolor{black} {#2}
      \]}
  \end{beamerboxesrounded}
}}

\AtBeginSection[] {
  \begin{frame}<beamer>\frametitle{Outline}
    \tableofcontents[currentsection,hideothersubsections]
  \end{frame}
}
%
%
%\AtBeginSubsection[] {
%  \begin{frame}<beamer>\frametitle{Outline}
%    \tableofcontents%[currentsection,currentsubsection]
%  \end{frame}
%}

%\usecolortheme[RGB={82,45,128}]{structure}
%\usecolortheme[RGB={162,80,22}]{structure}
\usecolortheme[RGB={128,0,0}]{structure}
\usetheme[secheader]{Boadilla}
%\usetheme[height=7mm]{Rochester}
%\usetheme{Copenhagen}
%\usetheme{Antibes}
%\usecolortheme{seahorse}
%\usecolortheme{crane}
%\usecolortheme{rose}
%\usefonttheme[onlylarge]{structurebold}
%\usefonttheme[onlymath]{serif}



\def\diag{{\rm diag}}


\def\E{\mathbb{E}}
\def\Prob{\mathbb{P}}
\def\argmin{{\rm argmin}}
\def\argmax{{\rm argmax}}
\def\Def{\stackrel{def}{=}}


\newtheorem{assumption}{Assumptions}
\newtheorem*{proposition}{Proposition}
\newtheorem*{remark}{Remark}



%\setbeamercolor{disc title}{bg=oiBlue!40!white!60,fg=blue}
\setbeamercolor{disc body}{bg= Regalia!20!white!80,fg= Regalia!80!black!90}

\setbeamercolor{clicker ungraded title}{bg=irishGreen!80!white!50,fg=irishGreen!30!black!90}
\setbeamercolor{clicker ungraded body}{bg=irishGreen!20!white!80,fg=irishGreen!30!black!90}

\setbeamercolor{clicker review title}{bg=gray!80!white!80,fg=oiBlue!80!black!90}
\setbeamercolor{clicker review body}{bg=gray!30!white!90,fg=oiBlue!80!black!90}

\setbeamercolor{code body}{bg=gray!20!white!80,fg=black}


% Custom commands
\newcommand{\degree}{\ensuremath{^\circ}}
\newcommand{\Note}[1]{
\rule{2.5cm}{0.25pt} \\ \textit{\scriptsize {\textcolor{rubineRed}{Note:} \textcolor{gray}{#1}}}}
\newcommand{\ct}[1]{
\vfill
{\tiny #1}}
\newcommand{\Remember}[1]{\textit{\scriptsize{\textcolor{orange}{Remember:} \textcolor{gray}{#1}}}}
\newcommand{\red}[1]{\textit{\textcolor{rubineRed}{#1}}}
\newcommand{\pink}[1]{\textit{\textcolor{rubineRed!90!white!50}{#1}}}
\newcommand{\green}[1]{\textit{\textcolor{irishGreen}{#1}}}
\newcommand{\webURL}[1]{\urlstyle{same}\textit{\textcolor{linkCol}{\url{#1}}} }
\newcommand{\webLink}[2]{\href{#1}{\textcolor{linkCol}{{#2}}}}
\newcommand{\mail}[1]{\href{mailto:#1}{\textit{\textcolor{linkCol}{#1}}}}
\newcommand{\hl}[1]{\textit{\textcolor{oiBlue}{#1}}}
\newcommand{\hlGr}[1]{\textit{\textcolor{lightGreen}{#1}}}
\newcommand{\mathhl}[1]{\textcolor{oiBlue}{\ensuremath{#1}}}
\newcommand{\ex}[1]{\textcolor{blue}{{{\small (#1)}}}}
\newcommand{\disc}[1]{
\begin{beamerboxesrounded}[shadow = true, lower = disc body, upper = disc title]{}
#1
\end{beamerboxesrounded}
}

\newcommand{\cl}[1]{
\begin{beamerboxesrounded}[shadow = true, lower = clicker ungraded body, upper = clicker ungraded title]{Question}
$\:$ \\
#1
\end{beamerboxesrounded}
}

\newcommand{\clR}[1]{
\begin{beamerboxesrounded}[shadow = true, lower = clicker review body, upper = clicker review title]{\red{Review question} }
$\:$ \\
#1
\end{beamerboxesrounded}
}

\newcommand{\formula}[2]{
\begin{beamerboxesrounded}[shadow = true, lower = white, upper = clicker review body]{#1}
#2
\end{beamerboxesrounded}
$\:$ \\
}

\newenvironment{twocol}[4]{
\begin{columns}[c]
\column{#1\textwidth}
#3
\column{#2\textwidth}
#4
\end{columns}
}


\newenvironment{slot}[2]{
\begin{array}{c}
\underline{#1} \\
#2
\end{array}
}

\newcommand{\pr}[1]{
\left( #1 \right)
}

\newcommand{\solnMult}[1]{
\item[] \vspace{-0.59cm}
\only<beamer| beamer:1>{\item #1}
\soln{\only<2->{\item \red{#1}}}
}

%\newcommand{\codechunk}[1]{
%\begin{beamerboxesrounded}[shadow = true, lower = code body]{}
%{\small #1}
%\end{beamerboxesrounded}
%}

% Change margin

\newenvironment{changemargin}[2]{%
\begin{list}{}{%
\setlength{\topsep}{0pt}%
\setlength{\leftmargin}{#1}%
\setlength{\rightmargin}{#2}%
\setlength{\listparindent}{\parindent}%
\setlength{\itemindent}{\parindent}%
\setlength{\parsep}{\parskip}%
}%
\item[]}{\end{list}}

% Footnote

\long\def\symbolfootnote[#1]#2{\begingroup%
\def\thefootnote{\fnsymbol{footnote}}\footnote[#1]{#2}\endgroup}

% Commands from the book
\newenvironment{data}[1]{\texttt{#1}}{}
\newenvironment{var}[1]{\texttt{#1}}{}
\newenvironment{resp}[1]{\texttt{#1}}{}






%%%%%%%%%%%%%%%%%%%%%%%%%%%%%%%%%%%%%%%%%%%%%%%%%%%%%%%%%%%%%%%%%%%%%%%%%%%%%%%%%%%%%%%%%%%%%%%

\title[Chapter 7 part 2]{Chapter 7 part 2}
\subtitle{Properties of Expectations}

%%%%%%%%%%%%%%%%%%%%%%%%%%%%%%%%%%%%%%%%%%%%%%%%%%%%%%%%%%%%%%%%%%%%%%%%%%%%%%%%%%%%%%%%%%%%%%%


\author[Jingchen (Monika) Hu] % (optional, use only with lots of authors)
{Jingchen (Monika) Hu}
% - Give the names in the same order as the appear in the paper.
% - Use the \inst{?} command only if the authors have different
%   affiliation.

\institute[Vassar] % (optional, but mostly needed)
{Vassar College}
% - Use the \inst command only if there are several affiliations.
% - Keep it simple, no one is interested in your street address.

\date[MATH 241] % (optional, should be abbreviation of conference name)
{MATH 241}
% - Either use conference name or its abbreviation.
% - Not really informative to the audience, more for people (including
%   yourself) who are reading the slides online

\subject{MATH 241}
% This is only inserted into the PDF information catalog. Can be left
% out.



% If you wish to uncover everything in a step-wise fashion, uncomment
% the following command:

%\beamerdefaultoverlayspecification{<+->}


\begin{document}

%\begin{frame}%[plain]
%\includegraphics[width = \textwidth]{figures/2015DukeNCAA}
%\end{frame}

%{ % all template changes are local to this group.
%\addtocounter{framenumber}{-1}
%    \setbeamertemplate{navigation symbols}{}
%    \begin{frame}[plain]
%        \begin{tikzpicture}[remember picture,overlay]
%            \node[at=(current page.center)] {
%                \includegraphics[width=1.25\paperwidth]{figures/2015DukeNCAA}
%            };
%        \end{tikzpicture}
%     \end{frame}
%}

%%%%%%%%%%%%%%%%%%%%%

% Title Page

\begin{frame}%[plain]
\titlepage
\end{frame}


%%%%%%%%%%%%%%%%%%%%%%
%\addtocounter{framenumber}{-1}
%
%\begin{frame}\frametitle{Annoucement}
%
%\begin{itemize}
%%\item HW7: \red{due now!}
%\item HW8: \red{due Tuesday, Nov 20th}
%
%
%\vspace{0.5cm}
%\item Course evaluation open tomorrow.
%  \begin{itemize}
%  \item Activate between 11/14 - 12/3.
%  \item If response rate $\geq 80\%$, drop the lowest quiz.
%  \end{itemize}
%
%\vspace{0.5cm}
%\item Next quiz: Tuesday, Nov 18th\\
%  \begin{itemize}
%  \item Topic: function of a continuous RV, i.e., find $f_Y(y)$ where $Y = g(X)$.
%  %\item To prepare: do homework questions in Chapter 5: 37, 39, 40, TE29
%  \end{itemize}
%
%
%%\item Midterm: Tuesday, Feb 25th
%%\begin{itemize}
%%\item Close book, in class exam (75 min)
%%\item ONE page cheat sheet {\bf made by yourself} (A4 size)
%%\item Calculators are allowed, but not cell phones, tablets or laptops
%%\end{itemize}
%
%\end{itemize}
%
%
%\end{frame}
%

%
%%%%%%%%%%%%%%%%%%%%%%
%\begin{frame}{Outline}
%%\tableofcontents[hideallsubsections,pausections]
%\tableofcontents[hideallsubsections]
%\end{frame}
%
%%%%%%%%%%%%%%%%%%%%%%%%%%%%%%%%%%%%%%%%%%%
%\begin{frame}\frametitle{Recap}
%
%
%\end{frame}





%%%%%%%%%%%%%%%%%%%%%
%\begin{frame}{Outline}
%%\tableofcontents[hideallsubsections,pausections]
%\tableofcontents[hideallsubsections]
%\end{frame}



%%%%%%%%%%%%%%%%%%%%%%%%%%%%%%%%%%%%%%%%%%
\section{Moment generating functions}
%%%%%%%%%%%%%%%%%%%%%%%%%%%%%%%%%%%%%%%%%%

\begin{frame}\frametitle{Moment generating functions}
\begin{block}{Definition}
The moment generating function $M(t)$ of the random variable $X$ is defined for all real values of $t$ by
\[ M(t) = E[e^{tX}] = \begin{cases}
    \sum_x e^{tx}p(x) & \text{if $X$ is discrete with pmf $p(x)$} \\
    \int_{-\infty}^{\infty}e^{tx}f(x)dx        & \text{if $X$ is continuous with pdf $f(x)$}
\end{cases}
 \]
\end{block}
\pause
{\small{
All the moments of $X$ can be obtained by successively differentiating $M(t)$ and then evaluating the result at $t = 0$\pause
$$M'(t) = \frac{d}{dt}E[e^{tX}] \pause= E[\frac{d}{dt}(e^{tX})] \pause= E[Xe^{tX}], \pause M'(0) = E[X]$$\pause
$$M''(t) = \frac{d}{dt}E[Xe^{tX}] \pause= E[\frac{d}{dt}(Xe^{tX})] \pause= E[X^2e^{tX}], \pause M''(0) = E[X^2]$$
\pause
In general
$$M^n(t) = E[X^ne^{tX}], M^n(0) = E[X^n], n \geq 1$$
}}
\end{frame}

\begin{frame}
\disc{If $X$ is a binomial random variable with parameters $n$ and $p$. Use the moment generating functions to obtain its mean and variance.}
\pause
{\small{
\begin{itemize}
\item The moment generating function is
\begin{eqnarray*}
M(t) = E[e^{tX}] &=& \sum_{k=0}^{n}e^{tk}{n \choose k}p^k (1-p)^{n-k} \\
&=& \sum_{k=0}^{n}{n \choose k} (pe^t)^k (1-p)^{n-k} \\
&=&(pe^t + 1 - p)^n
\end{eqnarray*}
\pause
\item Taking differentiations
$$M'(t) = n(pe^t + 1 - p)^{n-1} pe^t$$
\pause
$$M''(t) = n(n-1)(pe^t + 1 - p)^{n-2}(pe^t)^2 + n(pe^t + 1 - p)^{n-1}pe^t$$
\end{itemize}
}}
\end{frame}

\begin{frame}
$$M'(t) = n(pe^t + 1 - p)^{n-1} pe^t$$
$$M''(t) = n(n-1)(pe^t + 1 - p)^{n-2}(pe^t)^2 + n(pe^t + 1 - p)^{n-1}pe^t$$

Therefore
$$E[X] = M'(0) = np$$

\pause

\begin{eqnarray*}
Var(X) = E[X^2] - (E[X])^2 &=& M''(0) - (M'(0))^2 \\
&=& n(n-1)p^2 + np - (np)^2 \\
&=& np(1-p)
\end{eqnarray*}

\end{frame}

\begin{frame}{Properties of MGFs}
\begin{itemize}
\item Property 1: The MGF of the sum of independent random variables equals to the product of the individual MGFs.
$$M_{X+Y}(t) = E[e^{t(X+Y)}] = E[e^{tX}e^{tY}] = E[e^{tX}]E[e^{tY}] = M_X(t)M_Y(t)$$
\pause
\item Property 2: The MGF uniquely determines the distribution. Refer to the tables of MGF of some discrete and continuous distributions. {\color{red} Textbook pages 339 and 340 for lists of MGFs of distributions.}
\end{itemize}
\end{frame}


\begin{frame}{Joint moment generating functions}
\begin{block}{Definition}
For any $n$ random variables $X_1, ..., X_n$, the joint MGF $M(t_1, ..., t_n)$ is defined for all real values of $t_1, ..., t_n$ by
$$M(t_1, ..., t_n) = E[e^{t_1X_1+...t_nX_n}]$$
\end{block}%
\pause
{\small{
\begin{itemize}
\item The individual MGF can be obtained from $M(t_1, ..., t_n)$ by letting all but one of the $t_j$'s be 0
$$M_{X_i}(t) = E[e^{tX_i}] = M(0, ..., t, 0, ..., 0)$$
\pause
\item The joint MGF $M(t_1, ..., t_n)$ uniquely determines the joint distribution of $X_1, ..., X_n$
\pause
\item Then we have the $n$ random variables $X_1, ..., X_n$ are independent {\bf{if and only if}}
$$M(t_1, ..., t_n) = M_{X_1}(t_1)...M_{X_n}(t_n)$$
\end{itemize}
}}
\end{frame}
%
%%%%%%%%%%%%%%%%%%%%%%%%%%%%%%%%%%%%%%%%%%%
%\section{Additional properties of normal random variables}
%%%%%%%%%%%%%%%%%%%%%%%%%%%%%%%%%%%%%%%%%%%
%
%\begin{frame}\frametitle{Multivariate normal distribution}
%\begin{block}{Definition}
%{\small{
%Let $Z_1, ..., Z_n$ be a set of $n$ independent unit normal random variables. If, for some constants $a_{ij}$, $1 \leq i \leq m, 1 \leq j \leq n$, and $\mu_i, 1 \leq i \leq m$
%\begin{eqnarray*}
%X_1 &=& a_{11}Z_1 + ... + a_{1n}Z_n + \mu_1 \\
%X_2 &=& a_{21}Z_1 + ... + a_{2n}Z_n + \mu_2 \\
%\vdots && \vdots\\
%X_i &=& a_{i1}Z_1 + ... + a_{in}Z_n + \mu_i \\
%\vdots && \vdots\\
%X_m &=& a_{m1}Z_1 + ... + a_{mn}Z_n + \mu_m
%\end{eqnarray*}
%then the random variables $X_1, ..., X_m$ are said to have a multivariate normal distribution.
%}}
%\end{block}
%{\small{
%i.e. $X_1, ..., X_m$ are all linear combinations of a finite set of independent standard normal random variables.
%}}
%\end{frame}
%
%\begin{frame}\frametitle{Multivariate normal distribution}
%$$X_i = a_{i1}Z_1 + ... + a_{in}Z_n + \mu_i$$
%{\small{
%\begin{itemize}
%\item $E[X_i] = \mu_i, Var(X_i) = \sum_{j=1}^{n}a_{ij}^2$ for each $X_i$
%\pause
%\item Consider the joint MGF $M(t_1, ..., t_m) = E[e^{\{t_1X_1+...+t_mX_m\}}]$
%\begin{itemize}
%\item $\sum_{i=1}^{m}t_iX_i$ is a linear combination of $Z_1, ..., Z_n$
%\item The mean is
%$E[\sum_{i=1}^{m}t_iX_i] = \sum_{i=1}^{m}t_i\mu_i$
%\item The variance
%$$Var(\sum_{i=1}^{m}t_iX_i) = Cov(\sum_{i=1}^{m}t_iX_i, \sum_{j=1}^{m}t_jX_j) = \sum_{i=1}^{m}\sum_{j=1}^{m}t_it_jCov(X_i, X_j)$$
%\pause
%\item As $E[e^Y] = M_Y(t)|_{t=1} = e^{\mu + \sigma^2/2}$ if $Y \sim N(\mu, \sigma^2)$
%$$M(t_1,...,t_m) = e^{\{\sum_{i=1}^{m}t_i\mu_i + \frac{1}{2}\sum_{i=1}^{m}\sum_{j=1}^{m}t_it_jCov(X_i, X_j)\}}$$
%\end{itemize}
%\item {\color{red} The joint $X_1, ..., X_m$ is completely determined from $E[X_i]$ and $Cov(X_i, X_j), i, j = 1, ..., m$}
%\end{itemize}
%}}
%\end{frame}
%
%\begin{frame}
%\disc{Find $P(X < Y)$ for bivariate normal random variables $X$ and $Y$ having parameters: $X \sim N(\mu_x, \sigma_x^2), Y \sim N(\mu_y, \sigma_y^2), \rho = Corr(X, Y)$.}
%\pause
%\begin{itemize}
%\item The difference of $X$ and $Y$: $X - Y$
%\begin{itemize}
%\item $E[X - Y] = \mu_x - \mu_y$
%\item $Var(X - Y) = Var(X) + Var(-Y) + 2Cov(X, -Y) = \sigma_x^2 + \sigma_y^2 - 2\rho\sigma_x\sigma_y$
%\end{itemize}
%\pause
%\item Then we can obtain
%\begin{eqnarray*}
%P\{X < Y\} &=& P\{X - Y < 0\} \\ \pause
%&=&
%P\{\frac{X - Y - (\mu_x - \mu_y)}{\sqrt{\sigma_x^2 + \sigma_y^2 - 2\rho\sigma_x\sigma_y}} < \frac{- (\mu_x - \mu_y)}{\sqrt{\sigma_x^2 + \sigma_y^2 - 2\rho\sigma_x\sigma_y}}\}\\
%&=&
%\Phi\left(\frac{\mu_y - \mu_x}{\sqrt{\sigma_x^2 + \sigma_y^2 - 2\rho\sigma_x\sigma_y}}\right)
%\end{eqnarray*}
%\end{itemize}
%\end{frame}
\end{document}
